\section{Suplemento sobre a cohomologia de feixes}
\label{section:0.12}

\subsection{Cohomologia de feixes de módulos sobre espaços anelados}
\label{subsection:0.12.1}

\begin{env}[12.1.1]
\label{0.12.1.1}
\oldpage[0\textsubscript{III}]{49}
Seja $(X,\sh{O}_X)$ um espaço anelado.
Lembremos que, para todo $\sh{O}_X$-módulo $\sh{F}$, define-se a
cohomologia $\HH^\bullet(X,\sh{F})$, que é um funtor cohomológico universal
(T, 2.2) da categoria $\cat{C}(X)$ de
$\sh{O}_X$-módulos na categoria de grupos abelianos; é o funtor derivado
do funtor exato à esquerda
\[
  \sh{F}\longmapsto\Gamma(X,\sh{F}).
\]
O funtor $\HH^\bullet(X,\sh{F})$ é isomorfo
\oldpage[0\textsubscript{III}]{50}
à restrição à categoria $\cat{C}(X)$ do funtor cohomológico
definido do mesmo modo sobre a categoria de feixes de grupos abelianos
em $X$ (G, II, 7.2.1).
\end{env}

\begin{env}[12.1.2]
\label{0.12.1.2}
Ponhamos $A=\Gamma(X,\sh{O}_X)$.
Como todo elemento de $A$ define um endomorfismo do grupo abeliano
$\Gamma(X,\sh{F})$, define por funtorialidade um endomorfismo do
$\delta$-funtor $\HH^\bullet(X,\sh{F})$.
Estes endomorfismos dotam cada $\HH^p(X,\sh{F})$ de uma
estrutura de $A$-módulo, e o operador de conexão $\partial$ é
$A$-linear.
Além disso, para quaisquer dois inteiros não negativos $p,q$ e quaisquer dois
$\sh{O}_X$-módulos $\sh{F},\sh{G}$, existe um homomorfismo de
$A$-módulos, chamado \emph{produto cup},
\[
\label{0.12.1.2.1}
\tag{12.1.2.1}
  \HH^p(X,\sh{F})\otimes_A\HH^q(X,\sh{G})
  \longrightarrow
  \HH^{p+q}\bigl(X,\sh{F}\otimes_{\sh{O}_X}\sh{G}\bigr)
\]
(G, II, 6.6).
Estes homomorfismos fazem da soma direta $S$ dos
$\HH^p(X,\sh{O}_X)$, para $p\geq0$, uma $A$-álgebra graduada
anticomutativa, e da soma direta dos $\HH^p(X,\sh{F})$ um
$S$-módulo graduado.

Para todo recobrimento aberto $\mathfrak{U}$ de $X$, denotaremos sempre
por $\check{\mathrm{C}}^\bullet(\mathfrak{U},\sh{F})$ (contrariamente a
(G, II, 5.1)) o complexo de cocadeias alternadas do nervo de
$\mathfrak{U}$ com valores no sistema de coeficientes
$\Gamma(U_\alpha,\sh{F})$.
Está claro que
$\check{\mathrm{C}}^\bullet(\mathfrak{U},\sh{F})$ é um
$A$-módulo graduado, de modo que os grupos de cohomologia
$\check{\HH}^p(\mathfrak{U},\sh{F})$ deste complexo estão dotados de
uma estrutura de $A$-módulo.
Além disso, as aplicações canônicas
\[
  \check{\HH}^p(\mathfrak{U},\sh{F})
  \longrightarrow
  \HH^p(X,\sh{F})
\]
(G, II, 5.4) são $A$-lineares.
\end{env}

\begin{env}[12.1.3]
\label{0.12.1.3}
Seja $(X',\sh{O}_{X'})$ um segundo espaço anelado, e seja
$f=(\psi,\theta)$ um morfismo de $X'$ em $X$.
Ponhamos $A'=\Gamma(X',\sh{O}_{X'})$; o $\psi$-morfismo $\theta$
define canonicamente um homomorfismo de anéis $A\to A'$.
Seja $\sh{F}$ um $\sh{O}_X$-módulo e $\sh{F}'$ um
$\sh{O}_{X'}$-módulo.
Para todo $f$-morfismo $u:\sh{F}\to\sh{F}'$
\sref[0\textsubscript{I}]{0.4.4.1}, veremos que se pode definir,
para todo $p\geq0$, um dihomomorfismo
\[
\label{0.12.1.3.1}
\tag{12.1.3.1}
  u_p:\HH^p(X,\sh{F})\longrightarrow\HH^p(X',\sh{F}').
\]

Com efeito, como $\psi^*$ é exato na categoria de feixes de grupos
abelianos sobre $X$, o funtor
\[
  \sh{F}\longmapsto\HH^\bullet(X',\psi^*(\sh{F}))
\]
é um $\delta$-funtor sobre dita categoria, e sabe-se que existe um
homomorfismo canônico de $\delta$-funtores
\[
\label{0.12.1.3.2}
\tag{12.1.3.2}
  \HH^\bullet(X,\sh{F})
  \longrightarrow
  \HH^\bullet(X',\psi^*(\sh{F}))
\]
determinado de maneira única pela condição de que, em grau $0$, reduz-se
ao homomorfismo canônico
\[
  \Gamma(X,\sh{F})
  \longrightarrow
  \Gamma(X',\psi^*(\sh{F}))
\]
(T, 3.2.2).
Além disso, todo elemento de $A$ determina um endomorfismo $\mu$ de
$\Gamma(X,\sh{F})$ e um endomorfismo $\mu'$ de
$\Gamma(X',\psi^*(\sh{F}))$ tal que o diagrama
\[
\label{0.12.1.3.3}
\tag{12.1.3.3}
  \xymatrix{
    \Gamma(X,\sh{F})\ar[r]\ar[d]_\mu
      & \Gamma(X',\psi^*(\sh{F}))\ar[d]^{\mu'}\\
    \Gamma(X,\sh{F})\ar[r]
      & \Gamma(X',\psi^*(\sh{F}))
  }
\]
é comutativo.
\oldpage[0\textsubscript{III}]{51}
Pela propriedade de unicidade para prolongar morfismos de funtores
cohomológicos universais (T, 2.2), obtêm-se prolongações únicas de
$\mu$ e $\mu'$ à cohomologia que fazem comutativos os diagramas
análogos a \sref{0.12.1.3.3}.
Assim, \sref{0.12.1.3.2} é um homomorfismo de $A$-módulos.

Observemos agora que
\[
  f^*(\sh{F})
  =
  \psi^*(\sh{F})
  \otimes_{\psi^*(\sh{O}_X)}
  \sh{O}_{X'},
\]
e, portanto, que existe um dihomomorfismo canônico
\[
  \psi^*(\sh{F})\longrightarrow f^*(\sh{F})
\]
do $\psi^*(\sh{O}_X)$-módulo $\psi^*(\sh{F})$ no
$\sh{O}_{X'}$-módulo $f^*(\sh{F})$.
Por funtorialidade, obtém-se, pois, um dihomomorfismo funtorial
\[
\label{0.12.1.3.4}
\tag{12.1.3.4}
  \HH^p(X',\psi^*(\sh{F}))
  \longrightarrow
  \HH^p(X',f^*(\sh{F})),
\]
sendo $A$ e $A'$ os anéis correspondentes.
Compondo este dihomomorfismo com \sref{0.12.1.3.2}, obtém-se um
dihomomorfismo canônico, funtorial em $\sh{F}$,
\[
\label{0.12.1.3.5}
\tag{12.1.3.5}
  \theta_p:\HH^p(X,\sh{F})
  \longrightarrow
  \HH^p(X',f^*(\sh{F})).
\]

Finalmente, por funtorialidade, o homomorfismo
$u^\sharp:f^*(\sh{F})\to\sh{F}'$ dá um homomorfismo de
$A'$-módulos
\[
  \HH^p(X',f^*(\sh{F}))
  \longrightarrow
  \HH^p(X',\sh{F}'),
\]
que, composto com \sref{0.12.1.3.5}, dá
\sref{0.12.1.3.1}.

Seja $f'=(\psi',\theta'):X''\to X'$ um segundo morfismo de espaços
anelados, e seja $f''=f\circ f'$ o morfismo composto.
Levando em conta a compatibilidade do funtor $\psi'^*$ com os
produtos tensoriais \sref[0\textsubscript{I}]{0.4.3.3}, verifica-se
imediatamente que o composto do dihomomorfismo
\[
  \HH^p(X',f^*(\sh{F}))
  \longrightarrow
  \HH^p\bigl(X'',f'^*(f^*(\sh{F}))\bigr)
\]
com \sref{0.12.1.3.5} é o dihomomorfismo correspondente
\[
  \HH^p(X,\sh{F})
  \longrightarrow
  \HH^p(X'',f''^*(\sh{F})).
\]
\end{env}

\begin{env}[12.1.4]
\label{0.12.1.4}
Pode-se dar uma definição direta do homomorfismo
\sref{0.12.1.3.2} como segue.
Tomemos uma resolução injetiva
$\sh{L}^\bullet=(\sh{L}^i)$ de $\sh{F}$ constituída por feixes de
grupos abelianos sobre $X$.
Como o funtor $\psi^*$ é exato,
$\psi^*(\sh{L}^\bullet)$ é uma resolução de $\psi^*(\sh{F})$
constituída por feixes sobre $X'$.
Se $\sh{L}'^\bullet=(\sh{L}'^i)$ é uma resolução injetiva de
$\psi^*(\sh{F})$ na categoria de feixes de grupos abelianos sobre
$X'$, existe um morfismo
\[
  \psi^*(\sh{L}^\bullet)\longrightarrow\sh{L}'^\bullet
\]
de complexos de feixes de grupos abelianos, compatível com os
aumentos (M, V, 1.1.a), e determinado salvo homotopia.
Obtêm-se assim homomorfismos
\[
  \Gamma(X,\sh{L}^\bullet)
  \longrightarrow
  \Gamma(X',\psi^*(\sh{L}^\bullet))
  \longrightarrow
  \Gamma(X',\sh{L}'^\bullet)
\]
de complexos de grupos abelianos.
Seu composto, ao passar à cohomologia, dá um morfismo de
$\delta$-funtores
\[
  \HH^\bullet(X,\sh{F})
  \longrightarrow
  \HH^\bullet(X',\psi^*(\sh{F}));
\]
como coincide com \sref{0.12.1.3.2} em grau $0$, é idêntico a ele
(T, 2.2).

Consideremos agora um recobrimento aberto
$\mathfrak{U}=(U_\alpha)$ de $X$, e seja
$\mathfrak{U}'=(f^{-1}(U_\alpha))$ o recobrimento aberto por imagens inversas
de $X'$.
Para todo subconjunto aberto $V$ de $X$, os homomorfismos canônicos
\[
  \Gamma(V,\sh{F})
  \longrightarrow
  \Gamma(f^{-1}(V),f^*(\sh{F}))
\]
definem imediatamente (cf.\ G, II, 5.1) um homomorfismo de complexos
\[
  \check{\mathrm{C}}^\bullet(\mathfrak{U},\sh{F})
  \longrightarrow
  \check{\mathrm{C}}^\bullet(\mathfrak{U}',f^*(\sh{F})),
\]
e, portanto, homomorfismos canônicos
\[
\label{0.12.1.4.1}
\tag{12.1.4.1}
  \theta_p:
  \check{\HH}^p(\mathfrak{U},\sh{F})
  \longrightarrow
  \check{\HH}^p(\mathfrak{U}',f^*(\sh{F})).
\]

\oldpage[0\textsubscript{III}]{52}
Além disso, existem diagramas comutativos
\[
\label{0.12.1.4.2}
\tag{12.1.4.2}
  \xymatrix{
    \check{\HH}^p(\mathfrak{U},\sh{F})
      \ar[r]^-{\theta_p}\ar[d]
      & \check{\HH}^p(\mathfrak{U}',f^*(\sh{F}))\ar[d]\\
    \HH^p(X,\sh{F})
      \ar[r]_-{\theta_p}
      & \HH^p(X',f^*(\sh{F})),
  }
\]
onde as setas verticais são os homomorfismos canônicos de
(G, II, 5.2).
Para demonstrar a comutatividade de \sref{0.12.1.4.2}, consideremos o
complexo de feixes de cocadeias alternadas de $\sh{F}$ relativo a
$\mathfrak{U}$,
$\sh{C}^\bullet(\mathfrak{U},\sh{F})$, para o qual
\[
  \Gamma\bigl(X,\sh{C}^\bullet(\mathfrak{U},\sh{F})\bigr)
  =
  \check{\mathrm{C}}^\bullet(\mathfrak{U},\sh{F})
\]
(G, II, 5.2).
Os homomorfismos canônicos
\[
  \Gamma(V,\sh{F})
  \longrightarrow
  \Gamma(f^{-1}(V),\psi^*(\sh{F}))
\]
definem então um $\psi$-morfismo
\[
  \sh{C}^\bullet(\mathfrak{U},\sh{F})
  \longrightarrow
  \sh{C}^\bullet(\mathfrak{U}',\psi^*(\sh{F})),
\]
e, com a notação precedente, existe um diagrama comutativo
\[
  \xymatrix{
    \Gamma\bigl(X,\sh{C}^\bullet(\mathfrak{U},\sh{F})\bigr)
      \ar[r]\ar[d]
      & \Gamma\bigl(X',
          \sh{C}^\bullet(\mathfrak{U}',\psi^*(\sh{F}))\bigr)\ar[d]\\
    \Gamma(X,\sh{L}^\bullet)\ar[r]
      & \Gamma(X',\psi^*(\sh{L}^\bullet))\ar[r]
      & \Gamma(X',\sh{L}'^\bullet).
  }
\]
Ao passar à cohomologia, obtêm-se diagramas comutativos
\[
  \xymatrix{
    \check{\HH}^p(\mathfrak{U},\sh{F})\ar[r]\ar[d]
      & \check{\HH}^p(\mathfrak{U}',\psi^*(\sh{F}))\ar[d]\\
    \HH^p(X,\sh{F})\ar[r]
      & \HH^p(X',\psi^*(\sh{F})),
  }
\]
onde as setas verticais são os homomorfismos canônicos de
(G, II, 5.2).
Basta combinar estes diagramas com os diagramas comutativos
\[
  \xymatrix{
    \check{\HH}^p(\mathfrak{U}',\psi^*(\sh{F}))\ar[r]\ar[d]
      & \check{\HH}^p(\mathfrak{U}',f^*(\sh{F}))\ar[d]\\
    \HH^p(X',\psi^*(\sh{F}))\ar[r]
      & \HH^p(X',f^*(\sh{F})),
  }
\]
\oldpage[0\textsubscript{III}]{53}
que provêm do homomorfismo
$\psi^*(\sh{F})\to f^*(\sh{F})$ e da funtorialidade dos
homomorfismos canônicos de (G, II, 5.2), para obter a comutatividade
de \sref{0.12.1.4.2}.

Observemos que se
$A=\Gamma(X,\sh{O}_X)$ e $A'=\Gamma(X',\sh{O}_{X'})$, então
\sref{0.12.1.4.1} é um dihomomorfismo de módulos correspondente aos
anéis $A$ e $A'$.
Existe para \sref{0.12.1.4.1} uma propriedade de transitividade com respeito
a um composto de dois morfismos, análoga à de
\sref{0.12.1.3.5}.
Finalmente, nas definições precedentes, em lugar de uma resolução
injetiva $\sh{L}^\bullet$ de $\sh{F}$, poder-se-ia ter partido igualmente
de uma resolução tal que
\[
  \HH^p(X,\sh{L}^i)=0
  \qquad\text{para todo $i$ e todo $p>0$}
\]
(G, II, 4.7.1).
\end{env}

\begin{env}[12.1.5]
\label{0.12.1.5}
Sejam $\sh{F},\sh{G},\sh{H}$ três $\sh{O}_X$-módulos, e consideremos
um $\sh{O}_X$-homomorfismo
\[
  u:\sh{F}\otimes_{\sh{O}_X}\sh{G}\longrightarrow\sh{H},
\]
que dá, em cohomologia, os homomorfismos
\[
\label{0.12.1.5.1}
\tag{12.1.5.1}
  \HH^p(X,\sh{F})\otimes_A\HH^q(X,\sh{G})
  \longrightarrow
  \HH^{p+q}(X,\sh{H})
\]
deduzidos do produto cup \sref{0.12.1.2.1}.
Demonstraremos que, sob as hipóteses e com a notação de
\sref{0.12.1.3}, os diagramas
\[
\label{0.12.1.5.2}
\tag{12.1.5.2}
  \xymatrix@C=10pt{
    \HH^p(X,\sh{F})\otimes_A\HH^q(X,\sh{G})
      \ar[r]\ar[d]
      & \HH^{p+q}(X,\sh{H})\ar[d]\\
    \HH^p(X',f^*(\sh{F}))\otimes_{A'}
      \HH^q(X',f^*(\sh{G}))
      \ar[r]
      & \HH^{p+q}(X',f^*(\sh{H}))
  }
\]
são comutativos, onde as setas verticais provêm dos
homomorfismos canônicos \sref{0.12.1.3.5}.

Lembremos para isso que \sref{0.12.1.5.1} pode ser obtido a partir das
resoluções canônicas (G, II, 4.3)
$\sh{L}^\bullet,\sh{M}^\bullet,\sh{N}^\bullet$ de
$\sh{F},\sh{G},\sh{H}$ respectivamente, que estão constituídas por
$\sh{O}_X$-módulos, e da aplicação linear
\[
  \sh{L}^\bullet\otimes_{\sh{O}_X}\sh{M}^\bullet
  \longrightarrow
  \sh{N}^\bullet
\]
de complexos de $\sh{O}_X$-módulos correspondente a $u$.
Esta aplicação proporciona um homomorfismo de complexos de $A$-módulos
\[
  \Gamma(X,\sh{L}^\bullet)\otimes_A\Gamma(X,\sh{M}^\bullet)
  \longrightarrow
  \Gamma(X,\sh{N}^\bullet),
\]
e, portanto, ao passar à cohomologia, homomorfismos
\[
  \HH^p\bigl(\Gamma(X,\sh{L}^\bullet)\bigr)
  \otimes_A
  \HH^q\bigl(\Gamma(X,\sh{M}^\bullet)\bigr)
  \longrightarrow
  \HH^{p+q}\bigl(\Gamma(X,\sh{N}^\bullet)\bigr)
\]
(G, II, 6.6).
Existe claramente um diagrama comutativo
\[
\label{0.12.1.5.3}
\tag{12.1.5.3}
  \xymatrix@C=8pt{
    \Gamma(X,\sh{L}^\bullet)\otimes_A
      \Gamma(X,\sh{M}^\bullet)
      \ar[r]\ar[d]
      & \Gamma(X,\sh{N}^\bullet)\ar[d]\\
    \Gamma(X',\psi^*(\sh{L}^\bullet))
      \otimes_{\Gamma(X',\psi^*(\sh{O}_X))}
      \Gamma(X',\psi^*(\sh{M}^\bullet))
      \ar[r]
      & \Gamma(X',\psi^*(\sh{N}^\bullet)).
  }
\]
\oldpage[0\textsubscript{III}]{54}
Ao passar à cohomologia, obtêm-se diagramas comutativos
\[
\label{0.12.1.5.4}
\tag{12.1.5.4}
  \xymatrix@C=7pt{
    \HH^p(X,\sh{F})\otimes_A\HH^q(X,\sh{G})
      \ar[r]\ar[d]
      & \HH^{p+q}(X,\sh{H})\ar[d]\\
    \HH^p\bigl(\Gamma(X',\psi^*(\sh{L}^\bullet))\bigr)
      \otimes_{\Gamma(X',\psi^*(\sh{O}_X))}
      \HH^q\bigl(\Gamma(X',\psi^*(\sh{M}^\bullet))\bigr)
      \ar[r]
      & \HH^{p+q}\bigl(\Gamma(X',\psi^*(\sh{N}^\bullet))\bigr).
  }
\]
Mas como
$\psi^*(\sh{L}^\bullet)$,
$\psi^*(\sh{M}^\bullet)$, e
$\psi^*(\sh{N}^\bullet)$
são resoluções de
$\psi^*(\sh{F})$,
$\psi^*(\sh{G})$, e
$\psi^*(\sh{H})$ respectivamente, existe um diagrama comutativo
(G, II, 6.6.1)
\[
\label{0.12.1.5.5}
\tag{12.1.5.5}
  \xymatrix@C=7pt{
    \HH^p\bigl(\Gamma(X',\psi^*(\sh{L}^\bullet))\bigr)
      \otimes_{\Gamma(X',\psi^*(\sh{O}_X))}
      \HH^q\bigl(\Gamma(X',\psi^*(\sh{M}^\bullet))\bigr)
      \ar[r]\ar[d]
      & \HH^{p+q}\bigl(\Gamma(X',\psi^*(\sh{N}^\bullet))\bigr)
        \ar[d]\\
    \HH^p(X',\psi^*(\sh{F}))
      \otimes_{\Gamma(X',\psi^*(\sh{O}_X))}
      \HH^q(X',\psi^*(\sh{G}))
      \ar[r]
      & \HH^{p+q}(X',\psi^*(\sh{H})).
  }
\]
Finalmente, por funtorialidade, existe um diagrama comutativo
\[
\label{0.12.1.5.6}
\tag{12.1.5.6}
  \xymatrix@C=7pt{
    \HH^p(X',\psi^*(\sh{F}))
      \otimes_{\Gamma(X',\psi^*(\sh{O}_X))}
      \HH^q(X',\psi^*(\sh{G}))
      \ar[r]\ar[d]
      & \HH^{p+q}(X',\psi^*(\sh{H}))\ar[d]\\
    \HH^p(X',f^*(\sh{F}))\otimes_{A'}
      \HH^q(X',f^*(\sh{G}))
      \ar[r]
      & \HH^{p+q}(X',f^*(\sh{H})).
  }
\]
Combinando os três diagramas \sref{0.12.1.5.4},
\sref{0.12.1.5.5} e \sref{0.12.1.5.6}, obtém-se o diagrama
comutativo desejado \sref{0.12.1.5.2}.
\end{env}

\begin{remark}[12.1.6]
\label{0.12.1.6}
\oldpage[0\textsubscript{III}]{55}
Com a notação de \sref{0.12.1.3}, suponhamos que existe um
diagrama comutativo
\[
\label{0.12.1.6.1}
\tag{12.1.6.1}
  \xymatrix{
    0\ar[r] & \sh{F}\ar[r]^r\ar[d]_u
      & \sh{G}\ar[r]^s\ar[d]_v
      & \sh{H}\ar[r]\ar[d]_w & 0\\
    0\ar[r] & \sh{F}'\ar[r]_{r'}
      & \sh{G}'\ar[r]_{s'}
      & \sh{H}'\ar[r] & 0,
  }
\]
onde $r,s$ são homomorfismos de $\sh{O}_X$-módulos,
$r',s'$ são homomorfismos de $\sh{O}_{X'}$-módulos,
$u,v,w$ são $f$-morfismos, e as linhas são exatas.
Segue-se que o diagrama
\[
\label{0.12.1.6.2}
\tag{12.1.6.2}
  \xymatrix@C=6pt{
    \cdots\ar[r]
      & \HH^p(X,\sh{F})\ar[r]\ar[d]_{u_p}
      & \HH^p(X,\sh{G})\ar[r]\ar[d]_{v_p}
      & \HH^p(X,\sh{H})\ar[r]^-\partial\ar[d]_{w_p}
      & \HH^{p+1}(X,\sh{F})\ar[r]\ar[d]_{u_{p+1}}
      & \cdots\\
    \cdots\ar[r]
      & \HH^p(X',\sh{F}')\ar[r]
      & \HH^p(X',\sh{G}')\ar[r]
      & \HH^p(X',\sh{H}')\ar[r]_-\partial
      & \HH^{p+1}(X',\sh{F}')\ar[r]
      & \cdots
  }
\]
é comutativo.
Com efeito, \sref{0.12.1.6.1} se fatora como
\[
  \xymatrix@R=14pt{
    0\ar[r] & \sh{F}\ar[r]\ar[d]
      & \sh{G}\ar[r]\ar[d]
      & \sh{H}\ar[r]\ar[d] & 0\\
    0\ar[r] & \psi^*(\sh{F})\ar[r]\ar[d]
      & \psi^*(\sh{G})\ar[r]\ar[d]
      & \psi^*(\sh{H})\ar[r]\ar[d] & 0\\
    0\ar[r] & \sh{F}'\ar[r]
      & \sh{G}'\ar[r]
      & \sh{H}'\ar[r] & 0,
  }
\]
onde a linha intermediária é exata
\sref[0\textsubscript{I}]{0.3.7.2}.
Basta então utilizar o fato de que \sref{0.12.1.3.2} é um
homomorfismo de $\delta$-funtores e que os
$\HH^p(X',\sh{F}')$ formam um $\delta$-funtor em $\sh{F}'$.
\end{remark}

\begin{env}[12.1.7]
\label{0.12.1.7}
Com as hipóteses e a notação de \sref{0.12.1.3}, consideremos agora o
caso
\[
  \sh{F}=f_*(\sh{F}')=\psi_*(\sh{F}').
\]
Veremos que o dihomomorfismo definido em \sref{0.12.1.3},
\[
  \HH^p(X,f_*(\sh{F}'))\longrightarrow\HH^p(X',\sh{F}'),
  \tag{12.1.7.1}
  \label{0.12.1.7.1}
\]
pode ser obtido (salvo um automorfismo de $\HH^p(X',\sh{F}')$) como um
homomorfismo de bordo de uma sucessão espectral para o funtor composto
\[
  \sh{F}'\leadsto\Gamma(X,\psi_*(\sh{F}'))
\]
(T, 2.4).
A descrição de \sref{0.12.1.7.1} dada em \sref{0.12.1.4} mostra
aqui que este homomorfismo pode ser obtido como segue.
Tomemos resoluções injetivas $\sh{L}^\bullet$ e
${\sh{L}'}^\bullet$ de $\psi_*(\sh{F}')$ e $\sh{F}'$,
respectivamente, e depois um homomorfismo de complexos
\[
  v:\psi^*(\sh{L}^\bullet)\longrightarrow{\sh{L}'}^\bullet
\]
``acima'' do homomorfismo canônico
$\psi^*(\psi_*(\sh{F}'))\to\sh{F}'$.
\oldpage[0\textsubscript{III}]{56}
Observemos em seguida que
\[
  \Gamma(X',{\sh{L}'}^\bullet)
  =\Gamma(X,\psi_*({\sh{L}'}^\bullet))
\]
e que o homomorfismo composto
\[
  \Gamma(X,\sh{L}^\bullet)
  \longrightarrow\Gamma(X',\psi^*(\sh{L}^\bullet))
  \xrightarrow{\Gamma(v)}\Gamma(X',{\sh{L}'}^\bullet)
\]
não é senão
\[
  \Gamma(v^\flat):
  \Gamma(X,\sh{L}^\bullet)
  \longrightarrow\Gamma(X,\psi_*({\sh{L}'}^\bullet))
  \tag{12.1.7.2}
  \label{0.12.1.7.2}
\]
\sref[0\textsubscript{I}]{0.3.7.1}; e \sref{0.12.1.7.1} se
obtém passando à cohomologia em \sref{0.12.1.7.2}.

Por outro lado, as sucessões espectrais para o funtor composto
\[
  \sh{F}'\leadsto\Gamma(X,\psi_*(\sh{F}'))
\]
são obtidas tomando uma resolução injetiva de Cartan--Eilenberg
\[
  \sh{M}^{\bullet\bullet}=(\sh{M}^{ij})
\]
do complexo $\psi_*({\sh{L}'}^\bullet)$ na categoria de
feixes de grupos abelianos sobre $X$.
As sucessões espectrais em questão são as do bicomplexo
$\Gamma(X,\sh{M}^{\bullet\bullet})$; são birregulares, pois
$\sh{M}^{ij}=0$ para $i<0$ ou $j<0$.
Ora, a primeira sucessão espectral deste bicomplexo degenera, pois
os feixes $\psi_*({\sh{L}'}^i)$ são flácidos (G, II, 3.1.1), e
portanto
\[
  \HH_\mathrm{II}^q
    \bigl(\Gamma(X,\sh{M}^{i,\bullet})\bigr)
  =\HH^q\bigl(\psi_*({\sh{L}'}^i)\bigr)=0
  \qquad(q>0)
\]
(G, II, 4.4.3).
Por conseguinte, existem homomorfismos de bordo bijetivos
\sref{0.11.1.6}
\[
  {}'E_2^{i0}
  =\HH^i\!\left(
      \HH_\mathrm{II}^0
        \bigl(\Gamma(X,\sh{M}^{\bullet\bullet})\bigr)
    \right)
  \longrightarrow
  \HH^i\bigl(\Gamma(X,\sh{M}^{\bullet\bullet})\bigr),
  \tag{12.1.7.3}
  \label{0.12.1.7.3}
\]
e sabemos por \sref{0.11.3.4} que este homomorfismo provém, ao passar
à cohomologia, do aumento
\[
  \Gamma(X,\psi_*({\sh{L}'}^\bullet))
  \longrightarrow\Gamma(X,\sh{M}^{\bullet\bullet}),
  \tag{12.1.7.4}
  \label{0.12.1.7.4}
\]
que, por sua vez, provém do aumento
\[
  \eta:\psi_*({\sh{L}'}^\bullet)
  \longrightarrow\sh{M}^{\bullet\bullet}.
\]

Por outro lado, para a segunda sucessão espectral, existem homomorfismos
de bordo
\[
  {}''E_2^{i0}
  =\HH^i\!\left(
      \HH_\mathrm{I}^0
        \bigl(\Gamma(X,\sh{M}^{\bullet\bullet})\bigr)
    \right)
  \longrightarrow
  \HH^i\bigl(\Gamma(X,\sh{M}^{\bullet\bullet})\bigr),
  \tag{12.1.7.5}
  \label{0.12.1.7.5}
\]
que provêm, por \sref{0.11.3.4}, ao passar à cohomologia, do
homomorfismo de complexos
\[
  Z_\mathrm{I}^0
    \bigl(\Gamma(X,\sh{M}^{\bullet\bullet})\bigr)
  \longrightarrow\Gamma(X,\sh{M}^{\bullet\bullet}).
\]
Como $\psi_*$ é exato à esquerda, a sucessão
\[
  0\longrightarrow\psi_*(\sh{F}')
  \longrightarrow\psi_*({\sh{L}'}^0)
  \longrightarrow\psi_*({\sh{L}'}^1)
\]
é exata.
Pela definição de uma resolução de Cartan--Eilenberg
\sref{0.11.4.2}, podemos tomar, portanto,
\[
  B_\mathrm{I}^0(\sh{M}^{\bullet\bullet})=0,
  \qquad
  Z_\mathrm{I}^0(\sh{M}^{\bullet\bullet})=\sh{L}^\bullet.
\]
Como o diagrama
\[
  \xymatrix@C=50pt@R=24pt{
    \sh{L}^0\ar[r]^{i^0}
      & \sh{M}^{00}\\
    \psi_*(\sh{F}')\ar[u]^\varepsilon
      \ar[r]_{\varepsilon'}\ar[ur]^{\varepsilon''}
      & \psi_*({\sh{L}'}^0)\ar[u]_{\eta^0}
  }
\]
é comutativo, a injeção de complexos
$i:\sh{L}^\bullet\to\sh{M}^{\bullet\bullet}$ é compatível com os
aumentos $\varepsilon$ e $\varepsilon''$.
Existem, pois, dois homomorfismos de complexos de $\sh{L}^\bullet$
em $\sh{M}^{\bullet\bullet}$,
\[
  \xymatrix@C=52pt@R=22pt{
    \sh{L}^\bullet\ar[r]^i\ar[dr]_{v^\flat}
      & \sh{M}^{\bullet\bullet}\\
    & \psi_*({\sh{L}'}^\bullet)\ar[u]_\eta
  }
\]
\oldpage[0\textsubscript{III}]{57}
compatíveis com os aumentos $\varepsilon$ e
$\varepsilon''$.
Como $\sh{L}^\bullet$ é uma resolução injetiva e
$\sh{M}^{\bullet\bullet}$ está constituído por feixes injetivos, segue-se
de (M, V, 1.1(a)) que estes dois homomorfismos são homotópicos.
O mesmo vale, portanto, para os dois homomorfismos correspondentes
\[
  \Gamma(X,\sh{L}^\bullet)
  \longrightarrow\Gamma(X,\sh{M}^{\bullet\bullet}),
\]
e, ao passar à cohomologia, dão por conseguinte o mesmo
homomorfismo.
Em outras palavras, temos demonstrado que o homomorfismo de bordo
\sref{0.12.1.7.5}, escrito
\[
  \HH^p(X,\psi_*(\sh{F}'))
  \longrightarrow\HH^p\bigl(\Gamma(X,\sh{M}^{\bullet\bullet})\bigr),
\]
é o composto de \sref{0.12.1.7.1} e
\sref{0.12.1.7.3}, estando este último escrito
\[
  \HH^p(X',\sh{F}')
  \longrightarrow\HH^p\bigl(\Gamma(X,\sh{M}^{\bullet\bullet})\bigr)
\]
e tendo-se visto que é um isomorfismo.
Isto demonstra a asserção.
\end{env}

\subsection{Imagens diretas superiores}
\label{subsection:0.12.2}

\begin{env}[12.2.1]
\label{0.12.2.1}
Sejam $(X,\sh{O}_X)$ e $(Y,\sh{O}_Y)$ dois espaços anelados, e seja
$f=(\psi,\theta)$ um morfismo de $X$ em $Y$.  Esse morfismo define o funtor imagem
direta
\[
  f_*:\cat{C}(X)\longrightarrow\cat{C}(Y),
\]
que é além disso idêntico à restrição a $\cat{C}(X)$ do
funtor $\psi_*$ definido sobre a categoria de feixes de grupos abelianos
em $X$.
Este último funtor é aditivo e exato à esquerda; e como todo feixe de
grupos abelianos sobre $X$ é isomorfo a um subfeixe de um feixe injetivo
de grupos abelianos, definem-se os funtores derivados à direita
\[
  \sh{F}\longmapsto\RR^p\psi_*(\sh{F})
\]
de $\psi_*$.  Os $\RR^p\psi_*(\sh{F})$ são feixes de grupos abelianos
sobre $Y$, e os $\RR^p\psi_*$ formam um funtor cohomológico universal
(T, 2.3).

Além disso, $\RR^p\psi_*(\sh{F})$ é o feixe associado ao
pré-feixe
\[
  V\longmapsto\HH^p(f^{-1}(V),\sh{F})
\]
(T, 3.7.2).
Suponhamos agora que $\sh{F}$ é um $\sh{O}_X$-módulo.
Então $\HH^p(f^{-1}(V),\sh{F})$ está dotado naturalmente da estrutura
de $\Gamma(f^{-1}(V),\sh{O}_X)$-módulo, e portanto da de
$\Gamma(V,\psi_*(\sh{O}_X))$-módulo; o homomorfismo
\[
  \theta:\sh{O}_Y\longrightarrow\psi_*(\sh{O}_X)
\]
dota-o então da estrutura de
$\Gamma(V,\sh{O}_Y)$-módulo.
Para as estruturas assim definidas, a restrição de um subconjunto aberto
$V$ a um subconjunto aberto $V'\subset V$ define um dihomomorfismo.  Isto
define, portanto, sobre todo $\RR^p\psi_*(\sh{F})$ a estrutura de um
$\sh{O}_Y$-módulo.  Denotamos este $\sh{O}_Y$-módulo por
$\RR^p f_*(\sh{F})$; assim, $\RR^p f_*$ fica definido como um funtor
aditivo de $\cat{C}(X)$ em $\cat{C}(Y)$.

Além disso, os $\RR^p f_*$ formam um $\delta$-funtor.  Com efeito, se
\[
  0\longrightarrow\sh{F}'\longrightarrow\sh{F}
  \longrightarrow\sh{F}''\longrightarrow0
\]
é uma sucessão exata de $\sh{O}_X$-módulos, a descrição precedente
dos $\RR^p\psi_*$ e da estrutura de $\sh{O}_Y$-módulo sobre
$\RR^p\psi_*(\sh{F})$ mostra imediatamente que o homomorfismo de conexão
\[
  \partial:\RR^p\psi_*(\sh{F}'')
  \longrightarrow\RR^{p+1}\psi_*(\sh{F}')
\]
é um homomorfismo de $\sh{O}_Y$-módulos.
Finalmente, os $\RR^p f_*$ identificam-se com os funtores derivados à
direita de $f_*$.  Todo $\sh{O}_X$-módulo admite uma resolução injetiva por
$\sh{O}_X$-módulos, e tal resolução está constituída por feixes flácidos
de grupos abelianos (G, II, 7.1); pode, portanto, ser utilizada para calcular
os $\RR^p\psi_*(\sh{F})$, pois $\RR^n\psi_*(\sh{G})=0$ para $n\geq1$
para todo feixe flácido $\sh{G}$ (T, 2.4.1, Observação~3, e o
 corolário da Proposição~\EGAUnavailableHyperref{0.3.3.2}{3.3.2}).
Segue-se que os $\RR^p f_*$ formam um funtor cohomológico universal
de $\cat{C}(X)$ em $\cat{C}(Y)$ (T, 2.3).
\end{env}

\begin{env}[12.2.2]
\label{0.12.2.2}
Sejam $\sh{F}$ e $\sh{G}$ dois $\sh{O}_X$-módulos.
Com a notação de \sref{0.12.2.1}, para todo subconjunto aberto $V$ de
$Y$ existe o homomorfismo de produto cup \sref{0.12.1.2.1}
\[
  \begin{aligned}
  &\HH^p(f^{-1}(V),\sh{F})
    \otimes_{\Gamma(f^{-1}(V),\sh{O}_X)}
    \HH^q(f^{-1}(V),\sh{G})\\
  &\hspace{55mm}\longrightarrow
    \HH^{p+q}\bigl(f^{-1}(V),
      \sh{F}\otimes_{\sh{O}_X}\sh{G}\bigr).
  \end{aligned}
\]
\oldpage[0\textsubscript{III}]{58}
Segue-se imediatamente da definição do produto cup
(G, II, 6.6) que estes homomorfismos comutam com a restrição de
$V$ a um subespaço aberto $V'$ de $V$.
Por outro lado, $\theta$ dá um homomorfismo de anéis
\[
  \Gamma(V,\sh{O}_Y)
  \longrightarrow\Gamma(V,\psi_*(\sh{O}_X))
  =\Gamma(f^{-1}(V),\sh{O}_X),
\]
e, portanto, um homomorfismo canônico de produtos tensoriais
\[
  \begin{aligned}
  &\HH^p(f^{-1}(V),\sh{F})
    \otimes_{\Gamma(V,\sh{O}_Y)}
    \HH^q(f^{-1}(V),\sh{G})\\
  &\quad\longrightarrow
    \HH^p(f^{-1}(V),\sh{F})
    \otimes_{\Gamma(f^{-1}(V),\sh{O}_X)}
    \HH^q(f^{-1}(V),\sh{G}),
  \end{aligned}
\]
que é igualmente compatível com a restrição de $V$ a $V'$.
Por composição, obtém-se assim um homomorfismo de
$\Gamma(V,\sh{O}_Y)$-módulos.  Para os feixes associados aos
pré-feixes em questão, define um homomorfismo canônico,
funtorial em $\sh{F}$ e $\sh{G}$,
\[
  \RR^p f_*(\sh{F})\otimes_{\sh{O}_Y}\RR^q f_*(\sh{G})
  \longrightarrow
  \RR^{p+q}f_*
    \bigl(\sh{F}\otimes_{\sh{O}_X}\sh{G}\bigr).
  \tag{12.2.2.1}
  \label{0.12.2.2.1}
\]
Para $p=q=0$, este homomorfismo reduz-se a
\sref[0\textsubscript{I}]{0.4.2.2.1}.
\end{env}

\begin{proposition}[12.2.3]
\label{0.12.2.3}
Para todo $\sh{O}_X$-módulo $\sh{F}$ e todo $\sh{O}_Y$-módulo
$\sh{L}$ localmente livre de posto finito, existem isomorfismos canônicos
funtoriais
\[
  \RR^p f_*(\sh{F})\otimes_{\sh{O}_Y}\sh{L}
  \simto
  \RR^p f_*
    \bigl(\sh{F}\otimes_{\sh{O}_X}f^*(\sh{L})\bigr).
  \tag{12.2.3.1}
  \label{0.12.2.3.1}
\]
\end{proposition}

\begin{proof}
O homomorfismo \sref{0.12.2.3.1} é obtido compondo o
seguinte caso particular de \sref{0.12.2.2.1},
\[
  \RR^p f_*(\sh{F})
    \otimes_{\sh{O}_Y}f_*\bigl(f^*(\sh{L})\bigr)
  \longrightarrow
  \RR^p f_*
    \bigl(\sh{F}\otimes_{\sh{O}_X}f^*(\sh{L})\bigr),
  \tag{12.2.3.2}
  \label{0.12.2.3.2}
\]
com o homomorfismo do primeiro membro de
\sref{0.12.2.3.1} ao de \sref{0.12.2.3.2} induzido pelo
homomorfismo canônico \sref[0\textsubscript{I}]{0.4.4.3.2}.
Para verificar que \sref{0.12.2.3.1} é um isomorfismo quando $\sh{L}$ é
localmente livre, pode reduzir-se imediatamente ao caso
$\sh{L}=\sh{O}_Y$, pois a questão é local sobre $Y$ e os
funtores em questão são aditivos em $\sh{L}$.
Nesse caso, em vista da definição de \sref{0.12.2.2.1}, a
proposição reduz-se a verificar que o homomorfismo correspondente
de pré-feixes é bijetivo; isto segue-se imediatamente de
$f^*(\sh{O}_Y)=\sh{O}_X$.
\end{proof}

\begin{env}[12.2.4]
\label{0.12.2.4}
Seja $(Z,\sh{O}_Z)$ um terceiro espaço anelado, e seja
$g:Y\to Z$ um morfismo de espaços anelados.
Sabemos (G, II, 7.1 e 3.1.1) que, para todo
$\sh{O}_X$-módulo injetivo $\sh{G}$, $f_*(\sh{G})$ é um feixe flácido de
grupos abelianos.  Por conseguinte, por \sref{0.12.2.1},
\[
  \RR^p g_*\bigl(f_*(\sh{G})\bigr)=0
  \qquad(p>0).
\]
Segue-se (T, 2.4.1) que a sucessão espectral de Leray para funtores
compostos se aplica ao composto $g_*\circ f_*$.  Existe uma
sucessão espectral birregular cujo termo limite é o funtor
$\RR^\bullet h_*$, onde $h=g\circ f$, e cujo termo $E_2$ é
\[
  E_2^{pq}=\RR^p g_*\bigl(\RR^q f_*(\sh{F})\bigr).
  \tag{12.2.4.1}
  \label{0.12.2.4.1}
\]
\end{env}

\begin{env}[12.2.5]
\label{0.12.2.5}
Sob as condições de \sref{0.12.2.4}, definiremos diretamente
homomorfismos canônicos de $\sh{O}_Z$-módulos
\[
  \RR^n g_*\bigl(f_*(\sh{F})\bigr)
  \longrightarrow\RR^n h_*(\sh{F})
  \tag{12.2.5.1}
  \label{0.12.2.5.1}
\]
e
\[
  \RR^n h_*(\sh{F})
  \longrightarrow g_*\bigl(\RR^n f_*(\sh{F})\bigr),
  \tag{12.2.5.2}
  \label{0.12.2.5.2}
\]
\oldpage[0\textsubscript{III}]{59}
que podem identificar-se com os ``homomorfismos de bordo'' da sucessão
espectral de Leray (cf.\ \sref{0.12.1.7}).
Basta trabalhar com os pré-feixes aos quais estão associados os feixes de
imagens diretas superiores \sref{0.12.2.1}.
Para isso, seja $W$ um subconjunto aberto qualquer de $Z$, e consideremos
sua imagem inversa $g^{-1}(W)$ em $Y$.  Existe um
dihomomorfismo canônico
\[
  \begin{aligned}
  \HH^n\bigl(g^{-1}(W),f_*(\sh{F})\bigr)
  &\longrightarrow
  \HH^n\bigl(f^{-1}(g^{-1}(W)),
    f^*(f_*(\sh{F}))\bigr),
  \end{aligned}
  \tag{12.2.5.3}
  \label{0.12.2.5.3}
\]
sendo os anéis correspondentes
$\Gamma(g^{-1}(W),\sh{O}_Y)$ e
$\Gamma(h^{-1}(W),\sh{O}_X)$.
Por outro lado, o homomorfismo canônico
\sref[0\textsubscript{I}]{0.4.4.3.3} dá, por funtorialidade,
homomorfismos canônicos
\[
  \HH^n\bigl(h^{-1}(W),f^*(f_*(\sh{F}))\bigr)
  \longrightarrow
  \HH^n\bigl(h^{-1}(W),\sh{F}\bigr),
  \tag{12.2.5.4}
  \label{0.12.2.5.4}
\]
que são homomorfismos de
$\Gamma(h^{-1}(W),\sh{O}_X)$-módulos.
Levando em conta o homomorfismo de anéis
\[
  \Gamma(W,\sh{O}_Z)
  \longrightarrow\Gamma(g^{-1}(W),\sh{O}_Y),
\]
vemos que, compondo \sref{0.12.2.5.3} com
\sref{0.12.2.5.4}, se obtém um homomorfismo de pré-feixes e, portanto, o
homomorfismo de feixes \sref{0.12.2.5.1}.

A definição de \sref{0.12.2.5.2} é ainda mais simples.
Por definição, $\RR^n h_*(\sh{F})$ está associado ao pré-feixe
\[
  W\longmapsto
  \HH^n\bigl(f^{-1}(g^{-1}(W)),\sh{F}\bigr),
\]
enquanto que $\RR^n f_*(\sh{F})$ está associado ao pré-feixe
\[
  V\longmapsto\HH^n(f^{-1}(V),\sh{F}).
\]
Existe, portanto, um homomorfismo canônico
\[
  \HH^n\bigl(f^{-1}(g^{-1}(W)),\sh{F}\bigr)
  \longrightarrow
  \Gamma\bigl(g^{-1}(W),\RR^n f_*(\sh{F})\bigr),
\]
e é imediato que estes homomorfismos definem um homomorfismo de
pré-feixes, que, por sua vez, define \sref{0.12.2.5.2}.
\end{env}

\begin{env}[12.2.6]
\label{0.12.2.6}
Sob as hipóteses de \sref{0.12.2.4}, sejam $\sh{F}$, $\sh{G}$
e $\sh{K}$ três $\sh{O}_X$-módulos, e seja
\[
  u:\sh{F}\otimes_{\sh{O}_X}\sh{G}\longrightarrow\sh{K}
\]
um $\sh{O}_X$-homomorfismo.
Então os diagramas seguintes são comutativos:
\[
  \xymatrix@C=10pt@R=24pt{
    \RR^p g_*(f_*(\sh{F}))
      \otimes_{\sh{O}_Z}\RR^q g_*(f_*(\sh{G}))
      \ar[r]\ar[d]
      & \RR^{p+q}g_*(f_*(\sh{K}))\ar[d]\\
    \RR^p h_*(\sh{F})
      \otimes_{\sh{O}_Z}\RR^q h_*(\sh{G})
      \ar[r]
      & \RR^{p+q}h_*(\sh{K}),
  }
  \tag{12.2.6.1}
  \label{0.12.2.6.1}
\]
e
\[
  \xymatrix@C=10pt@R=24pt{
    \RR^p h_*(\sh{F})
      \otimes_{\sh{O}_Z}\RR^q h_*(\sh{G})
      \ar[r]\ar[d]
      & \RR^{p+q}h_*(\sh{K})\ar[d]\\
    g_*(\RR^p f_*(\sh{F}))
      \otimes_{\sh{O}_Z}g_*(\RR^q f_*(\sh{G}))
      \ar[r]
      & g_*(\RR^{p+q}f_*(\sh{K})).
  }
  \tag{12.2.6.2}
  \label{0.12.2.6.2}
\]
\oldpage[0\textsubscript{III}]{60}
Aqui as setas horizontais provêm de \sref{0.12.2.2.1} (a última
combinada com \sref[0\textsubscript{I}]{0.4.2.2.1}), e as setas
verticais de \sref{0.12.2.5.1} e \sref{0.12.2.5.2}, respectivamente.

Com efeito, basta verificar a asserção para os homomorfismos de pré-feixes
correspondentes.  Voltando às definições de
\sref{0.12.2.2} e \sref{0.12.2.5}, reduz-se imediatamente, para
\sref{0.12.2.6.1}, aos diagramas comutativos
\sref{0.12.1.5.2}; a verificação é ainda mais simples para
\sref{0.12.2.6.2}.
\end{env}

\subsection{Suplementos sobre os funtores Ext de feixes}
\label{subsection:0.12.3}

\begin{env}[12.3.1]
\label{0.12.3.1}
Seja $(X,\sh{O}_X)$ um espaço anelado.
Não recordaremos a definição nem as propriedades principais dos
bifuntores
\[
  \Ext^p_{\sh{O}_X}(X;\sh{F},\sh{G}),
  \qquad
  \shExt^p_{\sh{O}_X}(\sh{F},\sh{G}),
\]
o primeiro com valores na categoria de
$\Gamma(X,\sh{O}_X)$-módulos e o segundo na categoria de
$\sh{O}_X$-módulos, nem a sucessão espectral birregular
$E(\sh{F},\sh{G})$ que os relaciona (T, 4.2 e G, II, 7.3).
\end{env}

\begin{env}[12.3.2]
\label{0.12.3.2}
Assim como em (M, XIV, 1), define-se a noção de uma
\emph{extensão} de um $\sh{O}_X$-módulo $\sh{F}$ por um
$\sh{O}_X$-módulo $\sh{G}$ e a lei de composição sobre as classes de
extensões equivalentes: os argumentos utilizados para os módulos
transportam-se imediatamente para uma categoria abeliana arbitrária.
A segunda demonstração de (M, XIV, 1.1), que somente utiliza a existência de
imersões em objetos injetivos, continua sendo, portanto, válida na
categoria de $\sh{O}_X$-módulos.  Mostra que
\[
  \Ext^1_{\sh{O}_X}(X;\sh{F},\sh{G})
\]
identifica-se canonicamente com o grupo abeliano de classes de
extensões de $\sh{F}$ por $\sh{G}$.
\end{env}

\begin{proposition}[12.3.3]
\label{0.12.3.3}
Seja $(X,\sh{O}_X)$ um espaço anelado cujo feixe de anéis
$\sh{O}_X$ é coerente.
Então, para todo par de $\sh{O}_X$-módulos coerentes $\sh{F}$ e
$\sh{G}$ e todo $p\geq0$,
\[
  \shExt^p_{\sh{O}_X}(\sh{F},\sh{G})
\]
é um $\sh{O}_X$-módulo coerente.
\end{proposition}

\begin{proof}
Os $\shExt^p_{\sh{O}_X}(\sh{F},\sh{G})$ formam um funtor cohomológico
contravariante em $\sh{F}$.
Como $\sh{F}$ é coerente, para todo $p$ e todo ponto $x\in X$
existem uma vizinhança aberta $U$ de $x$ em $X$ e uma sucessão exata
de $(\sh{O}_X|U)$-módulos
\[
  0\longrightarrow\sh{R}\longrightarrow\sh{L}_{p-1}
  \longrightarrow\cdots\longrightarrow\sh{L}_0
  \longrightarrow\sh{F}|U\longrightarrow0,
\]
onde cada $\sh{L}_i$ ($0\leq i\leq p-1$) é isomorfo a
$\sh{O}_X^{n_i}|U$, e $\sh{R}$ é coerente.
Isto decorre por indução sobre $p$ de
\sref[0\textsubscript{I}]{0.5.3.2} e
\sref[0\textsubscript{I}]{0.5.3.4}, pois $\sh{O}_X$ é coerente.

Além disso, para $p\geq1$,
\[
  \shExt^p_{\sh{O}_X|U}(\sh{L}|U,\sh{G}|U)=0
\]
para todo $\sh{O}_X$-módulo $\sh{L}$ tal que $\sh{L}|U$ seja
isomorfo a $\sh{O}_X^n|U$ (T, 4.2.3).
O argumento de (M, V, 7.2) se aplica, portanto, ao funtor cohomológico
contravariante
\[
  \sh{F}\longmapsto
  \shHom_{\sh{O}_X|U}(\sh{F}|U,\sh{G}|U)
\]
e dá uma sucessão exata
\[
  \begin{aligned}
  \shHom_{\sh{O}_X|U}(\sh{L}_{p-1},\sh{G}|U)
  &\longrightarrow
  \shHom_{\sh{O}_X|U}(\sh{R},\sh{G}|U)\\
  &\longrightarrow
  \shExt^p_{\sh{O}_X|U}(\sh{F}|U,\sh{G}|U)
  \longrightarrow0.
  \end{aligned}
\]
Os dois primeiros termos são $(\sh{O}_X|U)$-módulos coerentes
\sref[0\textsubscript{I}]{0.5.3.5}; portanto, também o é o terceiro
\sref[0\textsubscript{I}]{0.5.3.4}.
\end{proof}

\oldpage[0\textsubscript{III}]{61}
\begin{proposition}[12.3.4]
\label{0.12.3.4}
Seja $f:X\to Y$ um morfismo plano de espaços anelados, e sejam
$\sh{F}$ e $\sh{G}$ dois $\sh{O}_Y$-módulos.
\begin{enumerate}
  \item[\rm(i)] Existe um homomorfismo de bifuntores cohomológicos
  \[
    f^*\bigl(\shExt^\bullet_{\sh{O}_Y}(\sh{F},\sh{G})\bigr)
    \longrightarrow
    \shExt^\bullet_{\sh{O}_X}
      \bigl(f^*(\sh{F}),f^*(\sh{G})\bigr),
    \tag{12.3.4.1}
    \label{0.12.3.4.1}
  \]
  que, em grau $0$, se reduz ao homomorfismo canônico
  \sref[0\textsubscript{I}]{0.4.4.6}.

  \item[\rm(ii)] Existe um morfismo canônico de sucessões espectrais
  \[
    E(\sh{F},\sh{G})
    \longrightarrow E\bigl(f^*(\sh{F}),f^*(\sh{G})\bigr),
    \tag{12.3.4.2}
    \label{0.12.3.4.2}
  \]
  cuja aplicação sobre os termos $E_2$ é dada pelos homomorfismos
  \[
    \begin{aligned}
    \HH^p\bigl(Y,\shExt^q_{\sh{O}_Y}(\sh{F},\sh{G})\bigr)
    \longrightarrow
    \HH^p\bigl(X,\shExt^q_{\sh{O}_X}
      (f^*(\sh{F}),f^*(\sh{G}))\bigr),
    \end{aligned}
    \tag{12.3.4.3}
    \label{0.12.3.4.3}
  \]
  deduzidos de \sref{0.12.3.4.1} e \sref{0.12.1.3.1}.
\end{enumerate}
\end{proposition}

\begin{proof}
\begin{enumerate}
  \item[\rm(i)]
  Como $f^*$ é exato sobre a categoria de $\sh{O}_Y$-módulos
  \sref[0\textsubscript{I}]{0.6.7.2}, os funtores
  \[
    \sh{G}\longmapsto
      f^*\bigl(\shHom_{\sh{O}_Y}(\sh{F},\sh{G})\bigr),
    \qquad
    \sh{G}\longmapsto
      \shHom_{\sh{O}_X}\bigl(f^*(\sh{F}),f^*(\sh{G})\bigr)
  \]
  são exatos à esquerda.  O homomorfismo canônico
  \sref[0\textsubscript{I}]{0.4.4.6} induz, portanto, canonicamente um
  homomorfismo entre seus funtores derivados.

  Para calcular estes funtores derivados, tomemos uma resolução injetiva
  $\sh{L}^\bullet=(\sh{L}^i)$ de $\sh{G}$.  Têm-se então
  homomorfismos
  \[
    \begin{aligned}
    \sh{H}^p\!\left(
      f^*\bigl(\shHom_{\sh{O}_Y}
        (\sh{F},\sh{L}^\bullet)\bigr)\right)
    \longrightarrow
    \sh{H}^p\!\left(
      \shHom_{\sh{O}_X}
        (f^*(\sh{F}),f^*(\sh{L}^\bullet))\right)
    \end{aligned}
  \]
  entre os feixes de cohomologia de complexos.
  Pela exatidão de $f^*$, o primeiro membro é
  \[
    f^*\bigl(\shExt^p_{\sh{O}_Y}(\sh{F},\sh{G})\bigr).
  \]
  A exatidão de $f^*$ implica também que
  $f^*(\sh{L}^\bullet)$ é uma resolução de $f^*(\sh{G})$.
  Se ${\sh{L}'}^\bullet=({\sh{L}'}^i)$ é uma resolução injetiva
  de $f^*(\sh{G})$ na categoria de $\sh{O}_X$-módulos, existe um
  homomorfismo de complexos
  \[
    f^*(\sh{L}^\bullet)\longrightarrow{\sh{L}'}^\bullet,
  \]
  determinado salvo homotopia, e portanto um homomorfismo bem definido em
  cohomologia.  Seu composto com o homomorfismo precedente é
  \sref{0.12.3.4.1}.

  \item[\rm(ii)]
  Com a notação precedente, existe um homomorfismo de complexos de
  $\sh{O}_X$-módulos
  \[
    f^*\bigl(\shHom_{\sh{O}_Y}(\sh{F},\sh{L}^\bullet)\bigr)
    \longrightarrow
    \shHom_{\sh{O}_X}
      \bigl(f^*(\sh{F}),{\sh{L}'}^\bullet\bigr).
  \]
  Seja $\sh{M}^{\bullet\bullet}$ uma resolução injetiva de
  Cartan--Eilenberg de
  $\shHom_{\sh{O}_Y}(\sh{F},\sh{L}^\bullet)$ na categoria de
  $\sh{O}_Y$-módulos.  Como $f^*$ é exato,
  $f^*(\sh{M}^{\bullet\bullet})$ é uma resolução de Cartan--Eilenberg de
  $f^*(\shHom_{\sh{O}_Y}(\sh{F},\sh{L}^\bullet))$.
  Se ${\sh{M}'}^{\bullet\bullet}$ é uma resolução injetiva de
  Cartan--Eilenberg de
  $\shHom_{\sh{O}_X}(f^*(\sh{F}),{\sh{L}'}^\bullet)$, existe
  \sref{0.11.4.2} um homomorfismo, determinado salvo homotopia,
  \[
    f^*(\sh{M}^{\bullet\bullet})
    \longrightarrow{\sh{M}'}^{\bullet\bullet},
  \]
  compatível com o homomorfismo precedente; em outras palavras, existe
  um $f$-morfismo
  $\sh{M}^{\bullet\bullet}\to{\sh{M}'}^{\bullet\bullet}$ de
  bicomplexos de feixes.
  Induz um dihomomorfismo
  \[
    \Gamma(Y,\sh{M}^{\bullet\bullet})
    \longrightarrow\Gamma(X,{\sh{M}'}^{\bullet\bullet})
  \]
  de bicomplexos de módulos, determinado salvo homotopia, e portanto
  um morfismo bem definido de sucessões espectrais
  \sref{0.11.3.2}.  Este é precisamente o morfismo desejado
  \sref{0.12.3.4.2}, e a descrição
  \sref{0.12.3.4.3} decorre imediatamente das definições.
\end{enumerate}
\end{proof}

\begin{proposition}[12.3.5]
\label{0.12.3.5}
Sob as hipóteses de \sref{0.12.3.4}, suponhamos além disso que o
feixe de anéis $\sh{O}_Y$ é coerente.
Então, para todo $\sh{O}_Y$-módulo coerente $\sh{F}$, os
homomorfismos canônicos \sref{0.12.3.4.1} são bijetivos.
\end{proposition}

\begin{proof}
\oldpage[0\textsubscript{III}]{62}
Como a questão é local sobre $Y$, podemos supor que existe uma
sucessão exata
\[
  0\longrightarrow\sh{R}\longrightarrow\sh{O}_Y^n
  \longrightarrow\sh{F}\longrightarrow0.
\]
Então $\sh{R}$ é também um $\sh{O}_Y$-módulo coerente
\sref[0\textsubscript{I}]{0.5.3.4}.
Para demonstrar que os homomorfismos
\[
  f^*\bigl(\shExt^p_{\sh{O}_Y}(\sh{F},\sh{G})\bigr)
  \longrightarrow
  \shExt^p_{\sh{O}_X}
    \bigl(f^*(\sh{F}),f^*(\sh{G})\bigr)
\]
são bijetivos, raciocinamos por indução sobre $p$.
Para $p=0$, a asserção é
\sref[0\textsubscript{I}]{0.6.7.6.1}.
Existe um diagrama comutativo com linhas exatas
\[
{\tiny
  \xymatrix@C=2pt@R=24pt{
  f^*\!\left(\shExt^{p-1}_{\sh{O}_Y}(\sh{O}_Y^n,\sh{G})\right)
    \ar[r]\ar[d]
  & f^*\!\left(\shExt^{p-1}_{\sh{O}_Y}(\sh{R},\sh{G})\right)
    \ar[r]^\partial\ar[d]
  & f^*\!\left(\shExt^p_{\sh{O}_Y}(\sh{F},\sh{G})\right)
    \ar[r]\ar[d]
  & f^*\!\left(\shExt^p_{\sh{O}_Y}(\sh{O}_Y^n,\sh{G})\right)
    \ar[d]\\
  \shExt^{p-1}_{\sh{O}_X}(\sh{O}_X^n,f^*(\sh{G}))
    \ar[r]
  & \shExt^{p-1}_{\sh{O}_X}(f^*(\sh{R}),f^*(\sh{G}))
    \ar[r]^\partial
  & \shExt^p_{\sh{O}_X}(f^*(\sh{F}),f^*(\sh{G}))
    \ar[r]
  & \shExt^p_{\sh{O}_X}(\sh{O}_X^n,f^*(\sh{G})).
  }}
\]
Aqui temos utilizado $f^*(\sh{O}_Y)=\sh{O}_X$; como $f^*$ é exato,
ambas as linhas são exatas.
Além disso,
\[
  \shExt^p_{\sh{O}_Y}(\sh{O}_Y^n,\sh{G})=0,
  \qquad
  \shExt^p_{\sh{O}_X}(\sh{O}_X^n,f^*(\sh{G}))=0
  \qquad(p>0)
\]
(T, 4.2.3).
Pela hipótese de indução, as duas primeiras setas verticais do
diagrama precedente são isomorfismos, enquanto que os termos da
direita são nulos.  Segue-se que
\[
  f^*\bigl(\shExt^p_{\sh{O}_Y}(\sh{F},\sh{G})\bigr)
  \longrightarrow
  \shExt^p_{\sh{O}_X}
    \bigl(f^*(\sh{F}),f^*(\sh{G})\bigr)
\]
é um isomorfismo.
\end{proof}

\subsection{Hipercohomologia do funtor imagem direta}
\label{subsection:0.12.4}

\begin{env}[12.4.1]
\label{0.12.4.1}
Sejam $(X,\sh{O}_X)$ e $(Y,\sh{O}_Y)$ dois espaços anelados, e seja
$f:X\longrightarrow Y$ um morfismo de espaços anelados.
Pode-se formar a hipercohomologia de $f_*$ com respeito a um complexo
arbitrário de $\sh{O}_X$-módulos
$\sh{K}^\bullet=(\sh{K}^i)_{i\in\bb{Z}}$
\sref{0.11.4.4}, pois os limites indutivos filtrantes existem e são exatos
na categoria abeliana de $\sh{O}_Y$-módulos (T, 3.1.1).
Os $\sh{O}_Y$-módulos de hipercohomologia
$\RR^p f_*(\sh{K}^\bullet)$ serão também denotados por
$\sh{H}^p(f,\sh{K}^\bullet)$, ou simplesmente
$\sh{H}^p_f(\sh{K}^\bullet)$.
Lembremos que $\sh{H}^\bullet(f,\sh{K}^\bullet)$ é a cohomologia do
bicomplexo de $\sh{O}_Y$-módulos $f_*(\sh{L}^{\bullet\bullet})$,
onde $\sh{L}^{\bullet\bullet}$ é uma resolução injetiva de
Cartan--Eilenberg de $\sh{K}^\bullet$ na categoria de
$\sh{O}_X$-módulos.
É o termo limite de duas sucessões espectrais cujos termos $E_2$
são dados por
\begin{equation}
\label{0.12.4.1.1}
\tag{12.4.1.1}
  {}'E_2^{pq}
  =\sh{H}^p\bigl(\sh{H}^q(f,\sh{K}^\bullet)\bigr)
\end{equation}
e
\begin{equation}
\label{0.12.4.1.2}
\tag{12.4.1.2}
  {}''E_2^{pq}
  =\sh{H}^p\bigl(f,\sh{H}^q(\sh{K}^\bullet)\bigr)
  \quad
  \bigl(=\RR^p f_*(\sh{H}^q(\sh{K}^\bullet))\bigr).
\end{equation}
Nestas fórmulas adotamos a notação geral $T(A^\bullet)$
para o transformado de um complexo por um funtor \sref{0.11.2.1}, e, para
um $\sh{O}_X$-módulo $\sh{F}$, escrevemos
$\sh{H}^p(f,\sh{F})$ em lugar de $\RR^p f_*(\sh{F})$.
Lembremos também que a primeira sucessão espectral é sempre regular.
Ambas as sucessões espectrais são birregulares quando $\sh{K}^\bullet$ está
limitado inferiormente,
\oldpage[0\textsubscript{III}]{63}
ou quando existe um inteiro $m$ tal que todo $\sh{O}_X$-módulo
admite uma resolução flácida de comprimento $\leq m$ \sref{0.11.4.4}.
\end{env}

\begin{env}[12.4.2]
\label{0.12.4.2}
Denotaremos analogamente por $\bb{H}^\bullet(X,\sh{K}^\bullet)$ a
hipercohomologia do funtor $\Gamma$ com respeito a um complexo
$\sh{K}^\bullet$ de $\sh{O}_X$-módulos.
Assim, os $\bb{H}^p(X,\sh{K}^\bullet)$ são
$\Gamma(X,\sh{O}_X)$-módulos.
Além disso, $\bb{H}^\bullet(X,\sh{K}^\bullet)$ pode ser considerado um
caso particular de $\sh{H}^\bullet(f,\sh{K}^\bullet)$, onde $f$ é um
morfismo de $(X,\sh{O}_X)$ a um espaço anelado constituído por um ponto
dotado do anel $\Gamma(X,\sh{O}_X)$.

Para todo subconjunto aberto $V$ de $X$, escreveremos
$\bb{H}^\bullet(V,\sh{K}^\bullet)$ em lugar de
$\bb{H}^\bullet(V,\sh{K}^\bullet|V)$.
\end{env}

\begin{proposition}[12.4.3]
\label{0.12.4.3}
Para todo inteiro $p\in\bb{Z}$, o $\sh{O}_Y$-módulo
$\sh{H}^p(f,\sh{K}^\bullet)$ é canonicamente isomorfo ao feixe
associado ao pré-feixe
\[
  U\longmapsto
  \bb{H}^p\bigl(f^{-1}(U),\sh{K}^\bullet\bigr)
\]
sobre $Y$.
\end{proposition}

\begin{proof}
Com a notação de \sref{0.12.4.1}, o feixe de cohomologia
$\sh{H}^p(f_*(\sh{L}^{\bullet\bullet}))$ está associado ao
pré-feixe
\[
  U\longmapsto
  \HH^p\bigl(\Gamma(U,f_*(\sh{L}^{\bullet\bullet}))\bigr)
  =
  \HH^p\bigl(\Gamma(f^{-1}(U),\sh{L}^{\bullet\bullet})\bigr).
\]
Mas está claro que
$\sh{L}^{\bullet\bullet}|f^{-1}(U)$ é uma resolução injetiva de
Cartan--Eilenberg de $\sh{K}^\bullet|f^{-1}(U)$ (T, 3.1.3).
Por conseguinte,
\[
  \HH^p\bigl(\Gamma(f^{-1}(U),\sh{L}^{\bullet\bullet})\bigr)
  =
  \bb{H}^p\bigl(f^{-1}(U),\sh{K}^\bullet\bigr)
\]
por definição.
\end{proof}

\begin{proposition}[12.4.4]
\label{0.12.4.4}
A hipercohomologia $\sh{H}^\bullet(f,\sh{K}^\bullet)$ é um
funtor cohomológico de $\sh{K}^\bullet$ em cada um dos casos
seguintes:
\begin{enumerate}
  \item[\rm{a)}] $\sh{K}^\bullet$ varia na categoria de complexos
    limitados inferiormente;
  \item[\rm{b)}] existe um inteiro $m$ tal que todo
    $\sh{O}_X$-módulo admite uma resolução flácida de comprimento $\leq m$;
  \item[\rm{c)}] $X$ é um espaço noetheriano.
\end{enumerate}
\end{proposition}

\begin{proof}
Os casos~a) e~b) são casos particulares de \sref{0.11.5.4}.
Por outro lado, o caso~c) segue-se de \sref{0.11.5.2}, pois, neste
caso, o funtor $f_*$ comuta com os limites indutivos
(G, II, 3.10.1).
\end{proof}

\begin{env}[12.4.5]
\label{0.12.4.5}
Seja agora $\mathfrak{U}=(U_\alpha)$ um recobrimento aberto de $X$.
Para todo complexo de \emph{pré-feixes}
$\sh{K}^\bullet=(\sh{K}^j)$ sobre $X$, consideremos o bicomplexo
$\check{\mathrm{C}}^\bullet(\mathfrak{U},\sh{K}^\bullet)$ cuja
componente de bigrado $(i,j)$ é
$\check{\mathrm{C}}^i(\mathfrak{U},\sh{K}^j)$, o grupo de
$i$-cocadeias alternadas do nervo de $\mathfrak{U}$ com valores
em $\sh{K}^j$ (G, II, 5.1).
Chamamos a cohomologia deste bicomplexo de
\emph{hipercohomologia do recobrimento $\mathfrak{U}$ com
coeficientes em $\sh{K}^\bullet$}, e a denotamos por
\[
  \bb{H}^\bullet(\mathfrak{U},\sh{K}^\bullet)
  =
  \HH^\bullet
  \bigl(\check{\mathrm{C}}^\bullet
    (\mathfrak{U},\sh{K}^\bullet)\bigr).
\]
A sucessão espectral de Leray de um recobrimento
(T, 3.8.1 e G, II, 5.9.1) se generaliza à hipercohomologia como segue.
\end{env}

\begin{proposition}[12.4.6]
\label{0.12.4.6}
Seja $\sh{K}^\bullet=(\sh{K}^i)$ um complexo de
$\sh{O}_X$-módulos.
Existe um funtor espectral regular de $\sh{K}^\bullet$ com
termo limite $\bb{H}^\bullet(X,\sh{K}^\bullet)$ e termo $E_2$
\begin{equation}
\label{0.12.4.6.1}
\tag{12.4.6.1}
  E_2^{pq}
  =
  \CHH^p\bigl(\mathfrak{U},h^q(\sh{K}^\bullet)\bigr),
\end{equation}
onde $h^q(\sh{K}^\bullet)$ denota o complexo de pré-feixes
\[
  V\longmapsto\HH^q(V,\sh{K}^\bullet)
\]
sobre $X$.
Esta sucessão espectral é birregular se $\sh{K}^\bullet$ está limitado
inferiormente.
\end{proposition}

\begin{proof}
Seja $\sh{L}^{\bullet\bullet}$ uma resolução injetiva de
Cartan--Eilenberg de $\sh{K}^\bullet$, e consideremos o tricomplexo
\[
  \check{\mathrm{C}}^\bullet
  (\mathfrak{U},\sh{L}^{\bullet\bullet})
  =
  \bigl(\check{\mathrm{C}}^i
    (\mathfrak{U},\sh{L}^{jk})\bigr).
\]
Consideremos primeiro este tricomplexo como um bicomplexo nos graus $i$ e
$j+k$.
Como $i$ somente toma valores não negativos, a segunda sucessão espectral
deste bicomplexo é regular \sref{0.11.3.3} e degenerada, porque
\[
  \CHH^q(\mathfrak{U},\sh{L}^{jk})=0
  \qquad(q>0),
\]
sendo os $\sh{O}_X$-módulos $\sh{L}^{jk}$ feixes flácidos
\oldpage[0\textsubscript{III}]{64}
(G, II, 5.2.3).
Segue-se que existe um isomorfismo canônico
\[
  \HH^n\bigl(
    \check{\mathrm{C}}^\bullet
      (\mathfrak{U},\sh{L}^{\bullet\bullet})
  \bigr)
  \isoto
  \HH^n\bigl(\Gamma(X,\sh{L}^{\bullet\bullet})\bigr)
\]
\sref{0.11.1.6}, em virtude de (G, II, 5.2.2), e portanto, pela
definição \sref{0.12.4.2}, um isomorfismo
\[
  \HH^n\bigl(
    \check{\mathrm{C}}^\bullet
      (\mathfrak{U},\sh{L}^{\bullet\bullet})
  \bigr)
  \isoto
  \bb{H}^n(X,\sh{K}^\bullet).
\]

Consideremos agora
$\check{\mathrm{C}}^\bullet
(\mathfrak{U},\sh{L}^{\bullet\bullet})$
como um bicomplexo nos graus $i+j$ e $k$.
Como $k$ somente toma valores não negativos, a primeira sucessão espectral
deste bicomplexo é sempre regular.
É birregular se $\sh{L}^{jk}=0$ para $j<j_0$, isto é, quando
$\sh{K}^\bullet$ está limitado inferiormente \sref{0.11.3.3}.
Esta é a sucessão espectral desejada.
Com efeito, para todo $j$, $\sh{L}^{j,\bullet}$ é uma resolução
injetiva de $\sh{K}^j$; por conseguinte, quando $i$ varia,
\[
  \HH^q\bigl(
    \check{\mathrm{C}}^i
      (\mathfrak{U},\sh{L}^{j,\bullet})
  \bigr)
\]
é precisamente o complexo de cocadeias
$\check{\mathrm{C}}^i(\mathfrak{U},h^q(\sh{K}^j))$.
Isto demonstra a proposição.
\end{proof}

\begin{corollary}[12.4.7]
\label{0.12.4.7}
Se, para todo simplexo $\sigma$ do nervo de $\mathfrak{U}$ e
todo inteiro $i$, tem-se
\[
  \HH^q(U_\sigma,\sh{K}^i)=0
  \qquad(q>0),
\]
então existe um isomorfismo canônico
\begin{equation}
\label{0.12.4.7.1}
\tag{12.4.7.1}
  \bb{H}^\bullet(\mathfrak{U},\sh{K}^\bullet)
  \isoto
  \bb{H}^\bullet(X,\sh{K}^\bullet).
\end{equation}
\end{corollary}

\begin{proof}
A hipótese implica que
$\check{\mathrm{C}}^\bullet
(\mathfrak{U},h^q(\sh{K}^j))=0$ para $q>0$, e portanto que
$E_2^{pq}=0$ para $q>0$.
Como a sucessão espectral \sref{0.12.4.6.1} é degenerada e
regular, a conclusão segue-se da definição
\sref{0.12.4.5} de
$\bb{H}^\bullet(\mathfrak{U},\sh{K}^\bullet)$ e de
\sref{0.11.1.6}.
\end{proof}

\begin{env}[12.4.8]
\label{0.12.4.8}
Seja $(X',\sh{O}_{X'})$ um segundo espaço anelado, e seja
$f=(\psi,\theta)$ um morfismo de $X'$ em $X$.
Pelo mesmo método que em \sref{0.12.1.3} e \sref{0.12.1.4}, se
define um dihomomorfismo para a hipercohomologia de um complexo
$\sh{K}^\bullet$ de $\sh{O}_X$-módulos:
\begin{equation}
\label{0.12.4.8.1}
\tag{12.4.8.1}
  \bb{H}^p(X,\sh{K}^\bullet)
  \longrightarrow
  \bb{H}^p(X',f^*(\sh{K}^\bullet)).
\end{equation}
Partindo de uma resolução injetiva de Cartan--Eilenberg
$\sh{L}^{\bullet\bullet}$ de $\sh{K}^\bullet$, observemos que, como
$\psi^*$ é exato, $\psi^*(\sh{L}^{\bullet\bullet})$ é uma
resolução de Cartan--Eilenberg de $\psi^*(\sh{K}^\bullet)$ na
categoria de $\psi^*(\sh{O}_X)$-módulos.
Existe, portanto, um morfismo
\[
  \psi^*(\sh{L}^{\bullet\bullet})
  \longrightarrow
  {\sh{L}'}^{\bullet\bullet},
\]
onde ${\sh{L}'}^{\bullet\bullet}$ é uma resolução injetiva de
Cartan--Eilenberg de $\psi^*(\sh{K}^\bullet)$.
Induz um morfismo em cohomologia
\[
  \bb{H}^\bullet(X,\sh{K}^\bullet)
  \longrightarrow
  \bb{H}^\bullet
  (X',\psi^*(\sh{K}^\bullet)).
\]
Compondo-o com o morfismo induzido por funtorialidade a partir de
$\psi^*(\sh{K}^\bullet)\longrightarrow f^*(\sh{K}^\bullet)$ obtém-se
o morfismo desejado \sref{0.12.4.8.1}.

Partindo de \sref{0.12.4.8.1} e \sref{0.12.4.3}, e raciocinando como
em \sref{0.12.2.5}, podem-se então definir, para dois morfismos de
espaços anelados
\[
  f:X\longrightarrow Y,
  \qquad
  g:Y\longrightarrow Z,
\]
com $h=g\circ f$, homomorfismos para a hipercohomologia de um complexo
$\sh{K}^\bullet$ de $\sh{O}_X$-módulos:
\begin{equation}
\label{0.12.4.8.2}
\tag{12.4.8.2}
  \sh{H}^n(g,f_*(\sh{K}^\bullet))
  \longrightarrow
  \sh{H}^n(h,\sh{K}^\bullet),
\end{equation}
e
\begin{equation}
\label{0.12.4.8.3}
\tag{12.4.8.3}
  \sh{H}^n(h,\sh{K}^\bullet)
  \longrightarrow
  g_*\bigl(\sh{H}^n(f,\sh{K}^\bullet)\bigr).
\end{equation}
Deixamos ao leitor os detalhes das definições.
\end{env}
